\section{Why the base character values already lie in the small field}

% BEGIN K77F1
\paragraph{K77F1. Recovering a ramified local character value.}
Let $E$ be a number field, $\chi$ a continuous idele class character
in the source normalization, and $M\subset\mathbb C$ a field
containing its unramified uniformizer values and every
$\chi_\infty(a)$ for $a\in E^\times$. Then
$\chi_v(y)\in M^\times$ for every finite $v$ and $y\in E_v^\times$.

At an unramified place write $y=\pi_v^d u$ with $u$ a unit.
Its value is $\chi_v(\pi_v)^d$. For the other places use continuity.
A disk $|z-1|<1/2$ contains no nontrivial subgroup of
$\mathbb C^\times$. For $|z|\ne1$ this follows by taking positive
or negative powers. On the unit circle the powers of any element
different from one meet the complement of this arc (double its
absolute argument until it exceeds $\pi/3$, starting inside the
arc). The inverse image of the disk contains a neighborhood of
the identity. Consequently $\chi$ kills local unit groups away
from finitely many primes and kills a sufficiently deep principal
unit subgroup at each remaining prime.

Let $S$ consist of the ramified primes together with $v$, and choose
one depth at least as large as the required depths in $S$ and at
least one. Weak approximation gives $a\in E^\times$ near $y$ at
$v$ and near $1$ at each $w\in S\setminus\{v\}$. At $v$ use the
open neighborhood $y(1+\mathfrak p_v^N)$, so this also works for
negative valuations of $y$. The principal idele of $a$ has character
value one. All its factors in $S$ other than $v$ are one, hence
\[
 \chi_v(y)=\left(\chi_\infty(a)
       \prod_{w\notin S}\chi_w(\pi_w)^{\operatorname{ord}_w(a)}
                    \right)^{-1}\in M^\times.
\]
The product is finite by factorization of $(a)$. This argument
recovers the value of an arbitrary local input, including a ramified
unit.
% END K77F1

% BEGIN K77F2
\paragraph{K77F2. Finite ideles and the assembled character.}
If $x$ has infinite part one, both its finite part and its inverse
are integral at almost every place. Thus $x_w$ is a unit almost
everywhere. Outside the finite union of its nonunit places and the
ramified places all local character factors are one. The truncated
ideles, with the omitted components set to one, converge to $x$ in
the restricted product topology as the finite set increases.
Their character values are eventually the same finite product.
Continuity and K77F1 give $\chi(x)\in M^\times$.

The host uses a circle-valued map $c$ on its own adelic algebra,
with the source comparison $\mathrm{ChrEq}(c,\chi)$. For this
presentation use the specified identification of the finite algebra
with the ordinary finite adeles. The idempotent for $w$ inserts a
local factor and is zero in every other component, including the
infinite part. Approximate the components of the given finite idele
at its nonunit and conductor primes by elements of $E^\times$ to
the depths of the open kernel. At other places units contribute one.
Equivalently, the finite product formula for $\mathrm{ChrEq}$ reads
$c(x)$ as a product of the corresponding $\chi_w(a_w)$. K77F1
therefore gives the same field assertion for $c$ in this presentation.
% END K77F2

% BEGIN K77F3
\paragraph{K77F3. The image of a pinned core field is exactly prescribed.}
Fix stages $X$, the distinguished embedding $e:L\hookrightarrow
\mathbb C$ and the prescribed small field $K_s\subset L$.
For a lawful core $J$ with its native field pin put $K=J.Q.c.F$,
an actual subfield of $\mathbb C$. The pin gives a bijection on
the integral coefficient records $e_O$ and identifies each integral
$a\in\mathcal O_K$ in $\mathbb C$ with $\phi_O(e_O a)$.
The field $F_O$ generated by the $\phi_O$-values therefore contains
$K$, by taking fractions of integers. Conversely every generator
comes from an $a$ by surjectivity of $e_O$. This proves $K=F_O$.
On assembled stages $\phi_O$ is $e$ on $\mathcal O_{K_s}$.
Taking fractions once more gives $F_O=e(K_s)$.

The algebra map $j:K\to L$ associated to this pin satisfies
$e(j(a))=a$ on $\mathcal O_K$ by its defining integral identity,
hence on $K$. Injectivity of $e$ now gives $j(K)=K_s$ inside $L$.
In particular the image lies in the original $L_0$ and contains
the image of $E$ under the specified host embedding, since the
embedded host coefficients belong to $F_O$. The core's ordered
primes, rank and canonical parameter pin are unaffected by the
identification. The field calculation applies to any such lawful
pinned core. Existence of the initial core uses the original
canonical parameter and adelic core constructions.
% END K77F3

% BEGIN K77F4
\paragraph{K77F4. A single coefficient field before the level choice.}
Retain the lawful pinned core and a pinned column with base characters
$f$. Its small-field pin contains the join of their value fields.
Let $L'$ be any allowed base level character record with $L'.C=f$
at these stages, and fix a row $i$. The unramified uniformizer
values of $\mu_{0,i+1}$ lie in $F_O$ by the corresponding summand
of that join. The infinity exponent at each embedding is even:
it is $-2z_\tau$ with $z_\tau$ the integral endpoint of the
stored weight. For nonzero $b$ the square of $b/|b|$ equals
$b/\bar b$. Its even power is thus an integral power of a quotient
of conjugate host embeddings. The host $E$ is Galois over
$\mathbb Q$, so every embedding of $E$, including its complex
conjugate, has the same image as the distinguished embedding.
K77F3 puts that image in $F_O$. Multiplication over the
embeddings puts $\chi_\infty(a)$ in $F_O$ for all $a\in E^\times$.

K77F2 applied to the level character and its pinned idele class
character puts its value on every finite idele in $F_O$. K77F3
then gives, for the same $j$ and every $x$ of infinite part one,
\[
           \mu_{0,i+1}(x)=e(j(y))\qquad\text{for some }y\in K .
\]
The witness $y$ may depend on $x,i,L'$. The field, its embedding,
and the two marked primes of the core were fixed before $L'$.
This is the value condition used in the later construction of
the analytic seed.
% END K77F4
